\chapter{تکنیک‌های پیشرفته تعامل با صفحه‌کلید}

سیستم‌عامل یونیکس اغلب به مزاح به‌عنوان «سیستم‌عاملی برای شیفتگان تایپ کردن» توصیف می‌شود و وجود رابط خط فرمان (\en{CLI})، گواهی بر این مدعاست. با این حال، کاربران حرفه‌ای خط فرمان در واقعیت چندان مایل به تایپ طولانی‌مدت نیستند. به همین سبب است که دستورات بنیادین با نام‌هایی بسیار کوتاه نظیر \cmd{cp}، \cmd{ls}، \cmd{mv} و \cmd{rm} طراحی شده‌اند. هدف غایی در معماری خط فرمان، بیشینه‌سازی بهره‌وری یا به بیانی دیگر، «تنبلی سازنده» است؛ یعنی انجامِ حجیم‌ترین عملیات با کمترین تعداد ضربه‌کلید. هدف دیگر این طراحی، متمرکز نگه داشتن کاربر بر صفحه‌کلید و جلوگیری از جابه‌جایی میان صفحه‌کلید و ماوس است.

در این فصل، برخی از قابلیت‌های پوسته \en{Bash} را بررسی خواهیم کرد که فرآیند تعامل با صفحه‌کلید را به طرز چشم‌گیری تسریع کرده و بهره‌وری کاربر را افزایش می‌دهند.

عناوین و دستوراتی که در این فصل به آن‌ها می‌پردازیم عبارتند از:

\begin{itemize}
  \item \cmd{clear}: پاک‌سازی محیط بصری ترمینال
  \item \cmd{history}: بازخوانی و نمایش فهرست تاریخچه دستورات اجرا شده
\end{itemize}

\section{ویرایش در خط فرمان}

پوسته \en{Bash} جهت مدیریت فرآیند ویرایش در خط فرمان، از کتابخانه \en{Readline} بهره می‌برد. این کتابخانه شامل مجموعه‌ای از روال‌ها (\en{Routines}) است که برنامه‌ها را قادر می‌سازد ورودی‌های کاربر را با انعطاف و کارایی بالاتری پردازش کنند. پیش‌تر با برخی از قابلیت‌های ابتدایی این کتابخانه، نظیر استفاده از کلیدهای جهت‌نما برای پیمایش، آشنا شده‌ایم؛ با این حال، \en{Readline} ابزارهای به‌مراتب پیشرفته‌تری را فراهم می‌آورد که می‌توانند سرعت تعامل با خط فرمان را به طرز محسوسی ارتقا دهند.

در ادامه، به بررسی برخی از این قابلیت‌های کلیدی می‌پردازیم. به یاد داشته باشید که نیازی به حفظ کردن تمامی این موارد نیست؛ بلکه توصیه می‌شود آن دسته از میان‌برهایی را که با الگوی کاری شما سازگارتر هستند برگزینید و به مرور در فعالیت‌های روزانه خود به کار بگیرید.

\begin{note}
\textbf{نکته:} برخی از ترکیب‌های کلیدی که شامل کلید \en{Alt} هستند، ممکن است توسط رابط گرافیکی (\en{GUI}) محیط دسکتاپ شما جهت عملکردهای سیستمی رزرو شده باشند. برای اطمینان از عملکرد صحیح تمامی توالی‌های کلیدی و جلوگیری از تداخل‌های نرم‌افزاری، استفاده از یک کنسول مجازی (\en{Virtual Console}) برای تمرین این موارد توصیه می‌شود.
\end{note}

\subsection{پیمایش مکان‌نما}

پوسته \en{Bash} جهت تسریع در جابه‌جایی مکان‌نما در خط فرمان، از مجموعه‌ای از کلیدهای میان‌بر بهره می‌برد. این میان‌برها به شما اجازه می‌دهند بدون نیاز به استفاده مکرر از کلیدهای جهت‌نما، با سرعت و دقت بالاتری در طول یک دستور پیمایش کنید.

\begin{table}[htbp]
\centering
\caption{دستورات پیمایش مکان‌نما}
\label{tab:cursor-navigation}
\begin{tabularx}{\textwidth}{l X}
\toprule
\textbf{میان‌بر} & \textbf{عملکرد فنی} \\
\midrule
\cmd{Ctrl+a} & انتقال آنی مکان‌نما به ابتدای خط. \\
\addlinespace
\cmd{Ctrl+e} & انتقال آنی مکان‌نما به انتهای خط. \\
\addlinespace
\cmd{Ctrl+f} & پیش‌روی به اندازه یک نویسه (معادل کلید جهت‌نمای راست). \\
\addlinespace
\cmd{Ctrl+b} & بازگشت به اندازه یک نویسه (معادل کلید جهت‌نمای چپ). \\
\addlinespace
\cmd{Alt+f} & جهش به کلمه بعدی (حرکت کلمه به کلمه رو به جلو). \\
\addlinespace
\cmd{Alt+b} & جهش به کلمه قبلی (حرکت کلمه به کلمه رو به عقب). \\
\addlinespace
\cmd{Ctrl+l} & پاک‌سازی نمایشگر و انتقال مکان‌نما به سطر نخست (عملکردی مشابه دستور \cmd{clear}). \\
\bottomrule
\end{tabularx}
\end{table}

\subsection{ویرایش متن}

از آنجا که خطاهای نوشتاری در زمان درج دستورات اجتناب‌ناپذیرند، بهره‌گیری از ابزارهایی جهت اصلاح سریع و بهینه آن‌ها ضرورتی دوچندان می‌یابد. جدول ذیل، کلیدهای میان‌بر اختصاصی برای ویرایش نویسه‌ها در محیط خط فرمان را شرح می‌دهد:

\begin{table}[htbp]
\centering
\caption{دستورات ویرایش متن}
\label{tab:text-editing}
\begin{tabularx}{\textwidth}{l X}
\toprule
\textbf{میان‌بر} & \textbf{عملکرد فنی} \\
\midrule
\cmd{Ctrl+d} & حذف نویسه‌ای که مکان‌نما دقیقاً بر روی آن قرار دارد. \\
\addlinespace
\cmd{Ctrl+t} & جابه‌جایی (\en{Transposition}) نویسه فعلی با نویسه ماقبل آن. \\
\addlinespace
\cmd{Alt+t} & جابه‌جایی کلمه فعلی با کلمه ماقبل آن. \\
\addlinespace
\cmd{Alt+l} & تبدیل تمامی حروف از موقعیت فعلی مکان‌نما تا انتهای کلمه به حروف کوچک (\en{Lowercase}). \\
\addlinespace
\cmd{Alt+u} & تبدیل تمامی حروف از موقعیت فعلی مکان‌نما تا انتهای کلمه به حروف بزرگ (\en{Uppercase}). \\
\bottomrule
\end{tabularx}
\end{table}

\subsection{برش و چسباندن متن}

در ادبیات فنی کتابخانه \en{Readline}، از مفاهیم برش (\en{Cut}) و چسباندن (\en{Paste}) تحت عناوین اختصاصی \en{``Killing''} و \en{``Yanking''} یاد می‌شود. قطعات متنی که طی این فرآیند حذف می‌شوند، در یک فضای ذخیره‌سازی موقت به نام \en{Kill-ring} نگهداری می‌گردند.

جدول زیر، کلیدهای میان‌بر ناظر بر این عملیات و منطق رفتاری هر یک را تبیین می‌کند:

\begin{table}[htbp]
\centering
\caption{دستورات برش و چسباندن متن}
\label{tab:cut-paste}
\begin{tabularx}{\textwidth}{l X}
\toprule
\textbf{میان‌بر} & \textbf{عملکرد فنی} \\
\midrule
\cmd{Ctrl+k} & برش متن از موقعیت فعلی مکان‌نما تا انتهای سطر. \\
\addlinespace
\cmd{Ctrl+u} & برش متن از موقعیت فعلی مکان‌نما تا ابتدای سطر. \\
\addlinespace
\cmd{Alt+d} & برش متن از موقعیت فعلی مکان‌نما تا انتهای کلمه جاری. \\
\addlinespace
\cmd{Alt+Backspace} & برش متن از موقعیت فعلی مکان‌نما تا ابتدای کلمه جاری؛ در صورت قرارگیری در ابتدای کلمه، کلمه پیشین حذف می‌شود. \\
\addlinespace
\cmd{Ctrl+y} & درج مجدد (\en{Yank}) آخرین محتوای موجود در \en{Kill-ring} در محل استقرار مکان‌نما. \\
\bottomrule
\end{tabularx}
\end{table}

\section{مفهوم کلید متا (Meta Key)}

در صورت مراجعه به مستندات فنی کتابخانه \en{Readline}، مکرراً با اصطلاح \en{Meta key} مواجه خواهید شد. در پیکربندی صفحه‌کلیدهای امروزی، این کلید به کلید \en{Alt} نگاشت (\en{Map}) شده است؛ هرچند این انطباق در گذشته به این شکل نبوده است.

در دوران پیش از ظهور کامپیوترهای شخصی، کاربران از طریق دستگاه‌هایی موسوم به ترمینال با سیستم تعامل داشتند. ترمینال‌ها ابزارهای ارتباطی محدودی بودند که صرفاً از یک نمایشگر متنی و یک صفحه‌کلید تشکیل شده و از طریق کابل سریال به مین‌فریم‌ها یا کامپیوترهای مرکزی متصل می‌شدند. به دلیل تنوع برندها و تفاوت در ویژگی‌های سخت‌افزاری این ترمینال‌ها، توسعه‌دهندگان نرم‌افزار برای تضمین قابلیت حمل (\en{Portability}) برنامه‌های خود، ناچار بودند به کمترین وجه مشترک سخت‌افزاری، یعنی مجموعه کاراکترهای \en{ASCII}، بسنده کنند. از آنجا که وجود کلیدهای کنترلی اضافی در تمامی این دستگاه‌ها تضمین‌شده نبود، یک کلید مجازی تحت عنوان «متا» تعریف شد تا عملکردهای فراتر از استاندارد را مدیریت کند.

امروزه، اگرچه کلید \en{Alt} نقش کلید متا را در سخت‌افزارهای مدرن ایفا می‌کند، اما همچنان می‌توان با فشردن و رها کردن سریع کلید \en{Esc} به همان کارکرد دست یافت. این قابلیت در لینوکس برای حفظ سازگاری عقب‌رو (\en{Backward Compatibility}) با محیط‌های ترمینال قدیمی و کنسول‌های خاص همچنان فعال باقی مانده است.

\section{سازوکار تکمیل خودکار (Completion)}

یکی از قدرتمندترین قابلیت‌های پوسته، مکانیزم تکمیل خودکار (\en{Completion}) است که با فشردن کلید \en{Tab} فعال می‌شود. این ویژگی به کاربر اجازه می‌دهد تا نام دستورات، مسیرها و فایل‌ها را با سرعت بالا و کمترین میزان تایپ به اتمام برساند.

برای درک بهتر فرآیند تکمیل، دایرکتوری خانگی (\en{Home Directory}) با محتوای زیر را در نظر بگیرید:

\begin{latin}
\begin{lstlisting}
[me@linuxbox ~]$ ls
Desktop     ls-output.txt  Pictures  Templates  Videos
Documents   Music          Public
\end{lstlisting}
\end{latin}

اگر بخشی از یک دستور را تایپ کرده و سپس کلید \en{Tab} را بفشارید، پوسته به‌طور خودکار نام فایل را تکمیل می‌کند؛ مشروط بر اینکه تنها یک گزینه با آن نویسه شروع شود. برای مثال:

\begin{latin}
\begin{lstlisting}
[me@linuxbox ~]$ ls l
\end{lstlisting}
\end{latin}

با فشردن کلید \en{Tab}، خروجی به شکل زیر کامل می‌شود:

\begin{latin}
\begin{lstlisting}
[me@linuxbox ~]$ ls ls-output.txt
\end{lstlisting}
\end{latin}

اما در صورتی که نویسه وارد شده با بیش از یک ورودی مطابقت داشته باشد، فرآیند تکمیل متوقف می‌شود:

\begin{latin}
\begin{lstlisting}
[me@linuxbox ~]$ ls D
\end{lstlisting}
\end{latin}

پس از فشردن کلید \en{Tab}:

\begin{latin}
\begin{lstlisting}
[me@linuxbox ~]$ ls D
\end{lstlisting}
\end{latin}

در این وضعیت، پوسته تغییری در خط فرمان ایجاد نمی‌کند، زیرا هر دو دایرکتوری \file{Desktop} و \file{Documents} با حرف \en{D} آغاز می‌شوند. برای موفقیت در فرآیند تکمیل، باید تعداد حروف کافی جهت شناسایی منحصربه‌فرد (\en{Unique Identification}) ورودی را تایپ کنید:

\begin{latin}
\begin{lstlisting}
[me@linuxbox ~]$ ls Do
\end{lstlisting}
\end{latin}

اکنون با فشردن کلید \en{Tab}، پوسته نام کامل را درج می‌کند:

\begin{latin}
\begin{lstlisting}
[me@linuxbox ~]$ ls Documents
\end{lstlisting}
\end{latin}

هرچند تکمیلِ مسیر و نام فایل متداول‌ترین کاربرد این مکانیزم است، اما قابلیت‌های تکمیل خودکار در پوسته \en{Bash} به حوزه‌های متنوع دیگری نیز تسری می‌یابد:

\begin{itemize}
  \item \textbf{متغیرها (\en{Variables}):} چنانچه واژه با نویسه \cmd{\$} آغاز شود، پوسته نام متغیرهای محیطی و محلی را پیشنهاد می‌دهد.
  \item \textbf{نام کاربری (\en{Usernames}):} اگر کلمه با نویسه \cmd{\textasciitilde} شروع شود، فرآیند تکمیل بر روی نام کاربران تعریف‌شده در سیستم انجام می‌گیرد.
  \item \textbf{دستورات (\en{Commands}):} در صورتی که واژه در ابتدای خط فرمان (موقعیت دستور) قرار داشته باشد، نام تمامی دستوراتِ موجود در مسیرهای اجرایی تعریف‌شده در متغیر محیطی \cmd{PATH} پیشنهاد می‌شود.
  \item \textbf{نام میزبان‌ها (\en{Hostnames}):} اگر کلمه با نویسه \cmd{@} آغاز گردد، پوسته نام میزبان‌های فهرست‌شده در فایل \file{/etc/hosts} را برای تکمیل استخراج می‌کند.
\end{itemize}

این سازوکار هوشمند با به حداقل رساندن خطاهای نوشتاری و صرفه‌جویی زمان، کارایی کاربر در محیط متنی را به شکل قابل توجهی ارتقا می‌دهد.

علاوه بر کلید \en{Tab}، چندین ترکیب کلیدی در گروه‌های \en{Control} و \en{Meta} وجود دارند که فرآیند تکمیل خودکار را تسهیل می‌کنند:

\begin{table}[htbp]
\centering
\caption{دستورات تکمیل خودکار}
\label{tab:completion}
\begin{tabularx}{\textwidth}{l X}
\toprule
\textbf{میان‌بر} & \textbf{عملکرد فنی} \\
\midrule
\cmd{Alt+?} & فهرستی از تمامی گزینه‌های ممکن برای تکمیل را نمایش می‌دهد. در اکثر توزیع‌های لینوکس، فشردن مجدد کلید \en{Tab} (دوبار فشردن پیاپی) همین عملکرد را ایفا می‌کند که روشی متداول‌تر است. \\
\addlinespace
\cmd{Alt+*} & تمامی گزینه‌های انطباقیِ ممکن را به‌طور مستقیم در خط فرمان درج می‌کند. این قابلیت زمانی که قصد دارید با مجموعه‌ای از نتایج منطبق به‌صورت هم‌زمان کار کنید، بسیار کارآمد است. \\
\bottomrule
\end{tabularx}
\end{table}

تعداد زیادی از این توالی‌های کلیدی با کاربردهای خاص و جزئی‌تر نیز وجود دارند. جهت بررسی فهرست جامع این میان‌برها، می‌توانید به بخش \en{READLINE} در صفحه راهنمای پوسته (\cmd{man bash}) مراجعه نمایید.

\section{تکمیل خودکار برنامه‌پذیر (Programmable Completion)}

نسخه‌های مدرن \en{Bash} از قابلیتی پیشرفته تحت عنوان تکمیل خودکار برنامه‌پذیر (\en{Programmable Completion}) بهره می‌برند. این ویژگی به کاربر یا به‌طور معمول به توسعه‌دهندگان توزیع‌های لینوکس اجازه می‌دهد تا قوانین تکمیلی اختصاصی را برای دستورات خاص تعریف کنند. هدف از این قابلیت، افزودن لایه‌ای از هوشمندی به پوسته است؛ به‌گونه‌ای که فرآیند تکمیل فراتر از نام فایل‌ها رفته و شامل گزینه‌ها (\en{Options})، آرگومان‌ها و حتی انواع فایل‌های مرتبط با یک برنامه خاص شود.

برای نمونه، هنگامی که دستوری مانند \cmd{git} را تایپ کرده و کلید \en{Tab} را می‌فشارید، پوسته بر اساس محتوای این سیستمِ کنترلی، زیردستورهای مرتبط (نظیر \cmd{clone}، \cmd{commit} یا \cmd{pull}) را پیشنهاد می‌دهد.

این قوانین در قالب توابع پوسته (\en{Shell Functions}) پیاده‌سازی شده‌اند که در واقع مینی‌اسکریپت‌های اجرایی درون پوسته هستند؛ مبحثی که در فصل‌های آتی با جزئیات فنی بیشتری به آن خواهیم پرداخت.

چنانچه نسبت به ساختار این توابع کنجکاو هستید، می‌توانید با اجرای دستور \cmd{set | less} تمامی توابع تعریف‌شده در محیط فعلی پوسته خود را بازبینی کنید. با این حال، شایان ذکر است که این ویژگی به‌صورت پیش‌فرض در تمام توزیع‌های لینوکس فعال نیست و ممکن است نیاز به نصب بسته‌های تکمیلی (مانند \cmd{bash-completion}) داشته باشد.

\section{تاریخچه دستورات}

پوسته \en{Bash} به‌طور خودکار فهرستی از تمامی دستورات ورودی شما را ثبت و نگهداری می‌کند. این سوابق در فایلی تحت عنوان \file{.bash\_history} واقع در دایرکتوری خانگی (\en{Home Directory}) ذخیره می‌شوند. قابلیت \en{History}، به‌ویژه هنگامی که با ابزارهای ویرایش خط فرمان ترکیب شود، راهکاری بی‌بدیل برای کاهش حجم تایپ و تکرار دستورات فراهم می‌آورد. این ویژگی به شما اجازه می‌دهد دستورات پیشین را به‌سرعت فراخوانی کرده، در صورت نیاز ویرایش نموده و مجدداً اجرا کنید.

\subsection{پیمایش و جستجو در تاریخچه}

شما می‌توانید در هر لحظه، محتوای فهرست تاریخچه دستورات خود را با اجرای فرمان زیر بازبینی کنید:

\begin{latin}
\begin{lstlisting}
[me@linuxbox ~]$ history | less
\end{lstlisting}
\end{latin}

به‌صورت پیش‌فرض، اکثر توزیع‌های مدرن لینوکس، \en{Bash} را به‌گونه‌ای پیکربندی کرده‌اند که ۱۰۰۰ مورد از آخرین دستورات اجرا شده را ذخیره نماید. شایان ذکر است که این مقدار کاملاً قابل شخصی‌سازی است؛ موضوعی که در فصل ۱۱ به تفصیل به آن خواهیم پرداخت.

تصور کنید قصد دارید تمامی دستورات اجرا شده پیرامون مسیر \file{/usr/bin/} را در سوابق خود بیابید. برای این کار می‌توانید خروجی دستور \cmd{history} را به ابزار \cmd{grep} پایپ (\en{Pipe}) کنید:

\begin{latin}
\begin{lstlisting}
[me@linuxbox ~]$ history | grep /usr/bin
\end{lstlisting}
\end{latin}

فرض کنید در میان نتایج جستجو، دستور کاربردی زیر را با شناسه اختصاصی ۸۸ مشاهده می‌کنید:

\begin{latin}
\begin{lstlisting}
88 ls -l /usr/bin > ls-output.txt
\end{lstlisting}
\end{latin}

این عدد در واقع شماره ردیف دستور در حافظه تاریخچه است. برای اجرای مجدد این فرمان بدون نیاز به بازنویسی کامل آن، می‌توانید از قابلیت بسط تاریخچه (\en{History Expansion}) استفاده کنید. این مکانیزم به شما اجازه می‌دهد با کمترین تلاش، دستورات پیشین را فراخوانی و اجرا نمایید.

جهت اجرای مجدد دستور شماره ۸۸، کافی است علامت تعجب (\cmd{!}) را به همراه شماره ردیف مربوطه درج کنید:

\begin{latin}
\begin{lstlisting}
[me@linuxbox ~]$ !88
\end{lstlisting}
\end{latin}

با فشردن کلید \en{Enter}، پوسته \en{Bash} به‌طور خودکار عبارت \cmd{!88} را با محتوای دقیق خط ۸۸ جایگزین و بلافاصله آن را اجرا می‌کند. این مورد تنها یکی از الگوهای بسط تاریخچه است؛ در بخش آتی، سایر فرمت‌ها و روش‌های پیشرفته فراخوانی از تاریخچه را بررسی خواهیم کرد.

پوسته \en{Bash} با ارائه قابلیت جستجوی افزایشی (\en{Incremental Search})، امکان پیمایش سریع و تعاملی در سوابق دستورات را فراهم می‌کند. در این فرآیند، با وارد کردن هر نویسه، نتایج جستجو به‌صورت آنی دقیق‌تر می‌شوند. این جستجو ماهیت «معکوس» (\en{Reverse}) دارد؛ به این معنا که پیمایش از جدیدترین دستورات ثبت‌شده به سمت قدیمی‌ترین‌ها صورت می‌گیرد.

\begin{latin}
\begin{lstlisting}
[me@linuxbox ~]$
\end{lstlisting}
\end{latin}

جهت فعال‌سازی این قابلیت، کلیدهای ترکیبی \cmd{Ctrl+r} را فشار دهید. در این حالت، پوسته وارد وضعیت جستجو شده و اعلان (\en{Prompt}) جدیدی را نمایش می‌دهد که آماده دریافت الگوی ورودی شماست:

\begin{latin}
\begin{lstlisting}
(reverse-i-search)`':
\end{lstlisting}
\end{latin}

با شروع تایپ عبارت مورد نظر، \en{Bash} به‌طور خودکار نخستین دستوری را که با الگوی واردشده مطابقت دارد، بازخوانی کرده و نمایش می‌دهد.

\textbf{مثال کاربردی:} تصور کنید به دنبال بازخوانی دستوری هستید که عبارت \file{/usr/bin} در آن به کار رفته است. پس از فشردن کلیدهای ترکیبی \cmd{Ctrl+r}، عبارت \file{/usr/bin} را تایپ می‌کنید:

\begin{latin}
\begin{lstlisting}
(reverse-i-search)`/usr/bin': ls -l /usr/bin > ls-output.txt
\end{lstlisting}
\end{latin}

به محض اینکه پوسته انطباقی (\en{Match}) را در تاریخچه شناسایی کند، دستور کامل را نمایش می‌دهد. در این مرحله، گزینه‌های زیر در اختیار شماست:

\begin{itemize}
  \item \textbf{اجرای مستقیم:} فشردن کلید \en{Enter} منجر به اجرای آنی دستور نمایش‌داده‌شده می‌شود.
  \item \textbf{انتقال برای ویرایش:} با فشردن کلیدهای \cmd{Ctrl+j}، دستور از حالت جستجو به خط فرمان فعلی منتقل می‌شود تا آماده ویرایش یا تغییر باشد.
  \item \textbf{ادامه پیمایش:} جهت یافتن موارد انطباقی قدیمی‌تر (تطابق بعدی)، مجدداً کلیدهای \cmd{Ctrl+r} را فشار دهید.
  \item \textbf{لغو عملیات:} برای خروج از حالت جستجو و بازگشت به اعلان (\en{Prompt}) عادی بدون اجرای دستور، از کلیدهای \cmd{Ctrl+c} یا \cmd{Ctrl+g} استفاده کنید.
\end{itemize}

با فشردن \cmd{Ctrl+j} در مثال فوق، دستور به خط فرمان اصلی شما منتقل می‌گردد:

\begin{latin}
\begin{lstlisting}
[me@linuxbox ~]$ ls -l /usr/bin > ls-output.txt
\end{lstlisting}
\end{latin}

اکنون این امکان فراهم است که دستور را پیش از اجرا، اصلاح نمایید. این قابلیت به‌طور چشم‌گیری سرعت عملیاتی شما را در مدیریت فرمان‌های پیچیده افزایش می‌دهد.

\begin{table}[htbp]
\centering
\caption{پیمایش و مدیریت تاریخچه دستورات}
\label{tab:history-navigation}
\begin{tabularx}{\textwidth}{l X}
\toprule
\textbf{میان‌بر} & \textbf{عملکرد فنی} \\
\midrule
\cmd{Ctrl+p} & فراخوانی دستور قبلی در تاریخچه (معادل کلید جهت‌نمای بالا). \\
\addlinespace
\cmd{Ctrl+n} & فراخوانی دستور بعدی در تاریخچه (معادل کلید جهت‌نمای پایین). \\
\addlinespace
\cmd{Alt+<} & انتقال مستقیم به ابتدای فهرست تاریخچه (قدیمی‌ترین دستور ثبت‌شده). \\
\addlinespace
\cmd{Alt+>} & انتقال مستقیم به انتهای فهرست تاریخچه (خط فرمان فعلی). \\
\addlinespace
\cmd{Ctrl+r} & فعال‌سازی جستجوی افزایشی معکوس (\en{Incremental Reverse Search})؛ با وارد کردن هر کاراکتر، جستجو به‌صورت آنی در سوابق انجام می‌شود. \\
\addlinespace
\cmd{Alt+p} & اجرای جستجوی معکوس غیرافزایشی؛ ابتدا الگوی جستجو را وارد کرده و پس از فشردن \en{Enter}، اولین انطباق بازیابی می‌شود. \\
\addlinespace
\cmd{Alt+n} & اجرای جستجوی مستقیم (رو به جلو) غیرافزایشی. \\
\addlinespace
\cmd{Ctrl+o} & اجرای دستور فعلی از تاریخچه و بارگذاری خودکار دستور بعدی؛ این میان‌بر برای تکرار یک توالی (\en{Sequence}) از دستورات بسیار کارآمد است. \\
\bottomrule
\end{tabularx}
\end{table}

\subsection{بسط تاریخچه (History Expansion)}

پوسته \en{Bash} جهت تسهیل دسترسی سریع به دستورات ثبت‌شده در فهرست سوابق، مکانیزم ویژه‌ای از انواع بسط (\en{Expansion}) را با محوریت نویسه علامت تعجب (\cmd{!}) ارائه می‌دهد. همان‌طور که پیش‌تر اشاره شد، درج یک عدد بلافاصله پس از علامت \cmd{!} منجر به فراخوانی دستور متناظر با آن ردیف در تاریخچه می‌شود. با این حال، قابلیت‌های بسط تاریخچه بسیار فراتر از فراخوانی ساده بر اساس شماره خط است.

این مکانیزم به شما اجازه می‌دهد تا با استفاده از الگوهای متنی یا موقعیت دستورات، بخش‌های خاصی از سوابق اجرایی خود را در دستور فعلی بازخوانی و ادغام کنید. جدول ذیل، جزئیات فنی و کاربردی این الگوهای بسط را تبیین می‌کند:

\begin{table}[htbp]
\centering
\caption{عملگرهای بسط تاریخچه}
\label{tab:history-expansion}
\begin{tabularx}{\textwidth}{l X}
\toprule
\textbf{توالی} & \textbf{عملکرد فنی} \\
\midrule
\cmd{!!} & تکرار آخرین دستور: آخرین فرمان اجرا شده را مجدداً فراخوانی و اجرا می‌کند. اگرچه در محیط‌های تعاملی، استفاده از کلید جهت‌نمای بالا (\en{Up Arrow}) و سپس فشردن \en{Enter} متداول‌تر است، اما این عملگر در ترکیب با دستورات دیگر (مانند \cmd{sudo !!}) بسیار کارآمد است. \\
\addlinespace
\cmd{!number} & فراخوانی بر اساس شماره ردیف: دستوری را که با شماره خط مشخص در فهرست تاریخچه ثبت شده است، اجرا می‌کند. شماره ردیف‌ها را می‌توان از طریق خروجی دستور \cmd{history} استخراج کرد. \\
\addlinespace
\cmd{!string} & فراخوانی بر اساس پیشوند: آخرین دستوری را که با رشته متنی (\en{String}) مشخص‌شده آغاز می‌شود، در تاریخچه جستجو و اجرا می‌کند. \\
\addlinespace
\cmd{!?string} & جستجوی محتوایی: آخرین دستوری را که حاوی رشته متنی مشخص‌شده در هر بخشی از بدنه خود باشد، یافته و اجرا می‌کند. \\
\bottomrule
\end{tabularx}
\end{table}

توصیه می‌شود هنگام استفاده از الگوهای \cmd{!string} و \cmd{!?string} احتیاط لازم را به عمل آورید، مگر آنکه از محتوای دقیق دستورات پیشین در حافظه تاریخچه اطمینان کامل داشته باشید. جهت کاهش مخاطرات ناشی از اجرای ناخواسته دستورات اشتباه، می‌توانید از اصلاح‌گر \cmd{:p} در انتهای عبارت استفاده کنید. این عملگر باعث می‌شود پوسته به‌جای اجرای آنی، ابتدا خروجی حاصل از بسط را نمایش داده و آن را به انتهای فهرست تاریخچه منتقل کند.

به عنوان مثال:

\begin{latin}
\begin{lstlisting}
[me@linuxbox ~]$ !ls:p
ls -l /usr/bin > ls-output.txt
\end{lstlisting}
\end{latin}

همان‌طور که ملاحظه می‌کنید، دستور مورد نظر صرفاً چاپ شده و به مرحله اجرا در نیامده است. اکنون که از صحت عبارت اطمینان یافتید، با توجه به اینکه دستور در انتهای سوابق قرار گرفته است، می‌توانید به‌سادگی با فشردن کلید جهت‌نمای بالا و \en{Enter} یا استفاده از فراخوان \cmd{!!} آن را اجرا کنید.

شایان ذکر است که خودِ عبارت‌های بسط (مانند \cmd{!!}) در تاریخچه ثبت نمی‌شوند، بلکه نتیجه نهایی و معادل بسط‌یافته آن‌ها در سوابق درج می‌گردد.

مکانیزم بسط تاریخچه واجد قابلیت‌های پیچیده‌تر و جزئیات فنی متعددی است که جهت حفظ انسجام مطالب، از ذکر آن‌ها صرف‌نظر می‌کنیم. در صورت تمایل به مطالعه عمیق‌تر، می‌توانید به بخش \en{HISTORY EXPANSION} در صفحه راهنمای پوسته (\cmd{man bash}) مراجعه فرمایید.

\section{دستور script: ضبط نشست‌های پوسته}

علاوه بر قابلیت‌های بومی تاریخچه در \en{Bash}، اکثر توزیع‌های لینوکس از برنامه‌ای کاربردی به نام \cmd{script} بهره می‌برند که قادر است تمامی کنش‌ها و واکنش‌های یک نشست (\en{Session}) در پوسته را ضبط و در قالب یک فایل متنی ذخیره نماید. این ابزار در حوزه‌هایی نظیر مستندسازی فرآیندهای سیستمی، تولید محتوای آموزشی یا عیب‌یابی (\en{Debugging}) پیشرفته، منبعی بسیار ارزشمند تلقی می‌شود.

نحو (\en{Syntax}) کلی استفاده از این دستور به شرح زیر است:

\begin{latin}
\begin{lstlisting}
script [file]
\end{lstlisting}
\end{latin}

در این الگو، \file{[file]} نشان‌دهنده نام فایلی است که محتوای نشست در آن ذخیره خواهد شد. چنانچه نام فایلی تعیین نگردد، برنامه به‌صورت پیش‌فرض خروجی را در فایلی تحت عنوان \file{typescript} ذخیره می‌کند.

جهت بازبینی فهرست جامع گزینه‌ها و ویژگی‌های فنی این ابزار، می‌توانید به صفحه راهنمای اختصاصی آن مراجعه کنید:

\begin{latin}
\begin{lstlisting}
man script
\end{lstlisting}
\end{latin}

این برنامه راهکاری ساده، استاندارد و بسیار مؤثر برای ثبت دقیق و گام‌به‌گام تعاملات کاربر با محیط خط فرمان فراهم می‌آورد.

\section{جمع‌بندی}

در این فصل، با مجموعه‌ای از تکنیک‌های عملیاتی در محیط پوسته آشنا شدیم که با هدف کاهش بار نوشتاری و ارتقای سرعت عمل کاربران حرفه‌ای طراحی شده‌اند. با افزایش تجربه و تسلط شما بر محیط خط فرمان، بازخوانی مجدد این سرفصل می‌تواند ابعاد تازه‌ای از بهره‌وری را برایتان آشکار سازد. در مقطع کنونی، این ابزارها را به عنوان قابلیت‌هایی کلیدی و راهبردی مد نظر قرار دهید که قادرند کارایی شما را در تعامل با سیستم‌عامل به شکل چشم‌گیری ارتقا بخشند.

\section{منابع مطالعه تکمیلی}

دانشنامه ویکی‌پدیا مقاله مناسبی درباره ترمینال‌های رایانه‌ای (\en{computer terminals}) دارد:

\begin{latin}
\url{http://en.wikipedia.org/wiki/Computer_terminal}
\end{latin}